%=============================================================================!
% 3.12  THE ENERGY OF ATOMS                                                   !
%=============================================================================!
\section{The energy of atoms}
\label{sec:en-atoms}

Every result of this chapter was found with boxes, beads and carts. This
section carries them to atoms: it introduces the unit in which the
energies of atoms are measured, converts the kinetic energy into the
units of this book, and proves that Newton's laws conserve the total
energy of any number of atoms whose forces come from a potential energy.

\subsection*{The electronvolt}

The joule is far too large a unit for a single atom. The lithium atom of
\cref{sec:ne-atoms}, of mass $6.94\times\SI{1.66054e-27}{\kilogram} =
\SI{1.1524e-26}{\kilogram}$, moving at \SI{1390}{\metre\per\second} after
\SI{10}{\femto\second} under a force of \SI{1}{\electronvolt\per
\angstrom}, has a kinetic energy of $\frac12\times1.1524\times10^{-26}
\times1390^2 = \SI{1.113e-20}{\joule}$. Energies of atoms are therefore
measured in electronvolts, the unit already used for forces in
\cref{sec:ne-atoms}, with
\begin{equation*}
   \SI{1}{\electronvolt} = \SI{1.602176634e-19}{\joule}
\end{equation*}
exactly. The name comes from electricity, which this book does not need;
the size of the unit in joules is fixed by the definition of the SI units,
and only its size matters here. The lithium atom has $1.113\times10^{-20}/1.602\times
10^{-19} = \SI{0.0695}{\electronvolt}$, and \cref{sec:ne-atoms} found
that it had moved \SI{0.0695}{\angstrom}. The force of
\SI{1}{\electronvolt\per\angstrom} has done
$1\times0.0695 = \SI{0.0695}{\electronvolt}$ of work on it, equal to its
kinetic energy, as the work-energy theorem requires.

\subsection*{Kinetic energy in eV, \AA, fs and amu}

With masses in atomic mass units and velocities in
\si{\angstrom\per\femto\second}, the product $\frac12 mv^2$ comes out in
\si{\amu\angstrom\squared\per\femto\second\squared}, which is not the
electronvolt; Newton's law needed a conversion factor in
\cref{sec:ne-atoms} for the same reason. We convert one unit of each
kind. A mass of \SI{1}{\amu} is \SI{1.66053906892e-27}{\kilogram}, and a
velocity of \SI{1}{\angstrom\per\femto\second} is $10^{-10}/10^{-15} =
\SI{e5}{\metre\per\second}$, whose square is
\SI{e10}{\metre\squared\per\second\squared}. Their product is, in joules,
\begin{equation*}
   \SI{1}{\amu\angstrom\squared\per\femto\second\squared}
   = 1.66053906892\times10^{-27}\times10^{10}\,\si{\joule}
   = \SI{1.66053906892e-17}{\joule} ,
\end{equation*}
since $-27 + 10 = -17$. Dividing by the size of an electronvolt in joules,
\begin{equation*}
   \frac{1.66053906892\times10^{-17}}{1.602176634\times10^{-19}}
   = 1.036427\times10^{2} = 103.6427 ,
\end{equation*}
where $1.66053906892/1.602176634 = 1.036427$ and $10^{-17}/10^{-19} =
10^{-17-(-19)} = 10^{2}$. Therefore
\begin{equation}
   \SI{1}{\amu\angstrom\squared\per\femto\second\squared}
   = \SI{103.6427}{\electronvolt} ,
   \qquad
   K\,[\si{\electronvolt}] = \tfrac12\,m\,[\si{\amu}]\;
   v^2\,[\si{\angstrom\squared\per\femto\second\squared}]
   \times103.6427 .
   \label{eq:en-mv2}
\end{equation}
This factor and the factor of Newton's law, \num{9.648533e-3}, are each
other's inverse, $1/103.6427 = \num{9.648533e-3}$, and for a reason.
\Cref{eq:en-mv2} says that one electronvolt is $1/103.6427$ of an
\si{\amu\angstrom\squared\per\femto\second\squared}, so a force of
\SI{1}{\electronvolt\per\angstrom}, divided by a mass of \SI{1}{\amu},
is $(1/103.6427)\,\si{\amu\angstrom\squared\per\femto\second\squared}$
divided by \si{\angstrom\amu}, which, cancelling the \si{\amu} and one
\si{\angstrom}, is $(1/103.6427)\,\si{\angstrom\per\femto\second\squared}$:
both factors
express the electronvolt in atomic mass units, \aa ngstr\"oms and
femtoseconds. In \texttt{mdlab} the kinetic-energy factor is
\texttt{MV2\_TO\_EV}, computed from the CODATA values of the constants,
and the Newton factor \texttt{FORCE\_TO\_ACCEL} is defined as its inverse,
so the two can never disagree. The kinetic energy of a group of atoms is
then a few lines.

\begin{code}
def kinetic_energy(masses: ArrayLike, velocities: ArrayLike) -> float:
    masses = np.asarray(masses, dtype=float)
    velocities = np.asarray(velocities, dtype=float)
    speed_squared = np.sum(velocities**2, axis=-1)  # v · v, one per atom
    if masses.ndim > 0 and masses.shape != speed_squared.shape:
        raise ValueError("give one mass for each row of velocities")
    in_amu_a2_per_fs2 = 0.5 * np.sum(masses * speed_squared)
    return float(in_amu_a2_per_fs2) * units.MV2_TO_EV
\end{code}

The velocities arrive with one row per atom and one column per component,
as the forces did in \cref{sec:ne-atoms}. Squaring the array squares every
component, and adding along the last axis, \texttt{axis=-1}, adds the
three squared components within each row, which is $\vb{v}_i\cdot
\vb{v}_i$ for each atom $i$. The check applies when the masses come as a
list, for which \texttt{masses.ndim}, the number of dimensions of the
array, is 1: it stops the function with a message unless
\texttt{masses.shape}, the length of each dimension, matches that of the
squared speeds, one mass for each atom. A single number, which has no
dimensions, skips the check and serves as the mass of every atom. Multiplying the masses by the squared
speeds pairs each atom's mass with its own speed, the sum over atoms and
the factor one half give $\sum_i\frac12 m_i v_i^2$ in
\si{\amu\angstrom\squared\per\femto\second\squared}, and the last line
converts to electronvolts. The lithium atom, which moves along a line
and so has a velocity of one component, called as
\texttt{kinetic\_energy(6.94, [0.0139])}, returns
\SI{0.0695}{\electronvolt}.

\subsection*{The potential energy of many atoms}

For $N$ atoms the potential energy is a function of all their positions,
$U(\vb{r}^N)$, a function of $3N$ coordinates. Chapter~8 shows where such
functions come from; here we need only that the force on each atom is
minus the gradient of $U$ with respect to that atom's own coordinates,
\begin{equation}
   \vb{F}_i = -\grad_i U
   = -\Bigl(\frac{\partial U}{\partial x_i}, \frac{\partial U}{\partial
   y_i}, \frac{\partial U}{\partial z_i}\Bigr) ,
   \label{eq:en-force-atoms}
\end{equation}
where the subscript on $\grad_i$ says which atom's coordinates are
differentiated, all the others held fixed. Everything in
\cref{sec:en-paths,sec:en-gradient,sec:en-curl} carries over with $3N$
coordinates in place of three: a path moves all the atoms at once, from
one arrangement to another; the work done on them along it is the sum
over atoms of $\vb{F}_i\cdot\Delta\vb{r}_i$ added over its short steps,
and for these forces it is the drop in $U$ between the two
arrangements.

\begin{proposition}
\label{prop:en-energy-atoms}
Let $N$ atoms move according to Newton's second law, $m_i\,\dd\vb{v}_i/\dd
t = \vb{F}_i$, with forces $\vb{F}_i = -\grad_i U(\vb{r}^N)$ from a
potential energy that does not depend on time except through the
positions. Then the total energy
\begin{equation*}
   E = \sum_{i=1}^{N}\tfrac12 m_i v_i^2 + U(\vb{r}^N)
\end{equation*}
is constant along every motion.
\end{proposition}

\begin{proof}
By \cref{sec:en-varying}, applied to each atom, the kinetic energy changes
at the rate $\dd K/\dd t = \sum_i\vb{F}_i\cdot\vb{v}_i$, where $K =
\sum_i\frac12 m_i v_i^2$, the sum of the
powers of the forces; this is where the second law enters. The gradient
toolbox of \cref{sec:en-gradient}, with a step of $3N$ legs, one along
each coordinate, in place of two, gives, grouping the legs atom by atom,
\begin{equation*}
   \Delta U \approx \sum_{i=1}^{N}\Bigl(\frac{\partial U}{\partial x_i}
   \Delta x_i + \frac{\partial U}{\partial y_i}\Delta y_i
   + \frac{\partial U}{\partial z_i}\Delta z_i\Bigr)
   = \sum_{i=1}^{N}\grad_i U\cdot\Delta\vb{r}_i ,
\end{equation*}
and the chain rule along a path then reads
\begin{equation*}
   \frac{\dd U}{\dd t} = \sum_{i=1}^{N}\grad_i U\cdot\frac{\dd\vb{r}_i}{\dd
   t} = -\sum_{i=1}^{N}\vb{F}_i\cdot\vb{v}_i .
\end{equation*}
The two rates cancel, so $\dd E/\dd t = 0$.
\end{proof}

A potential energy that depends only on the distance between two atoms,
like the energy $\frac12 k\xi^2$ of the spring between the carts of
\cref{sec:en-bodies}, is of this form, and \cref{ex:en-pair} shows that
for it \cref{eq:en-force-atoms}
produces the equal and opposite forces of the third law by itself.

\Needspace{10\baselineskip}
\begin{bridge}
In Born--Oppenheimer molecular dynamics, $U(\vb{r}^N)$ is the
ground-state energy of the electrons with the nuclei held at $\vb{r}^N$,
plus the repulsion between the nuclei, and the forces are its negative
gradient. The Hellmann--Feynman theorem supplies the part that comes from
the explicit dependence of the Hamiltonian on the nuclear positions, and
the Pulay terms add the part that comes from basis functions moving with
the atoms. A force computed without the Pulay terms, or from an
incompletely converged density, is not exactly $-\grad_i U$ of the energy
reported with it, and \cref{sec:en-drift} shows what such a discrepancy
does to the total energy.
\end{bridge}

\begin{selfcheck}
An oxygen atom (\SI{15.999}{\amu}) moves at
\SI{0.01}{\angstrom\per\femto\second}. What is its kinetic energy in
electronvolts?
\end{selfcheck}
